% Beispieldokument für das Beamer-Thema NXT
\documentclass[aspectratio=169]{beamer}

\usetheme{NXT}
\usepackage{booktabs}
\usepackage[table]{xcolor}
\usepackage[ngerman]{babel}

% Farbmuster: gefülltes Rechteck + Farbname
\newcommand{\cs}[1]{%
  \colorbox{#1}{\kern1.5em\strut\kern1.5em}~{\footnotesize\ttfamily #1}%
}

% Metadaten
\title{NXT-Folien mit \LaTeX{}}
\subtitle{Beispiel für das Beamer-Thema NXT}
\author{Volker Reichenberger}
\date{\today}

\begin{document}

%%%%%%%%%%%%%%%%%%%%%%%%%%%%%%%%%%%%%%%%%%%%%%%%%%%%%%%%%%%%%%%%%%%%%%%%%%%%%%%
% Titelseite
\begin{frame}[plain]
  \titlepage
\end{frame}

%%%%%%%%%%%%%%%%%%%%%%%%%%%%%%%%%%%%%%%%%%%%%%%%%%%%%%%%%%%%%%%%%%%%%%%%%%%%%%%
\begin{frame}{Aufzählungen}

\begin{columns}[T]
  \begin{column}{0.48\textwidth}
    \textbf{Aufzählung (itemize)}
    \begin{itemize}
      \item Erster Punkt
      \item Zweiter Punkt mit \emph{Hervorhebung}
      \item Dritter Punkt
        \begin{itemize}
          \item Unterpunkt A
          \item Unterpunkt B
        \end{itemize}
    \end{itemize}
  \end{column}
  \begin{column}{0.48\textwidth}
    \textbf{Nummerierung (enumerate)}
    \begin{enumerate}
      \item Schritt eins
      \item Schritt zwei
      \item Schritt drei
    \end{enumerate}
  \end{column}
\end{columns}

\end{frame}

%%%%%%%%%%%%%%%%%%%%%%%%%%%%%%%%%%%%%%%%%%%%%%%%%%%%%%%%%%%%%%%%%%%%%%%%%%%%%%%
\begin{frame}{Textblöcke und Hervorhebungen}

\begin{block}{Hinweis}
  Mit der Umgebung \texttt{block} können wichtige Informationen
  hervorgehoben werden.
\end{block}

\bigskip

\begin{alertblock}{Warnung}
  Diese Umgebung eignet sich für Warnungen oder kritische Hinweise.
\end{alertblock}

\bigskip

\begin{exampleblock}{Beispiel}
  Und diese Umgebung für Beispiele oder Erläuterungen.
\end{exampleblock}

\end{frame}

%%%%%%%%%%%%%%%%%%%%%%%%%%%%%%%%%%%%%%%%%%%%%%%%%%%%%%%%%%%%%%%%%%%%%%%%%%%%%%%
\begin{frame}{Tabellen}

\begin{table}
  \centering
  \begin{tabular}{llr}
    \toprule
    \textbf{Methode} & \textbf{Beschreibung} & \textbf{Aufwand} \\
    \midrule
    Interviews       & Qualitative Befragung        & hoch   \\
    Fragebogen       & Quantitative Erhebung        & mittel \\
    Beobachtung      & Direkte Verhaltensanalyse    & hoch   \\
    Dokumentanalyse  & Auswertung vorhandener Daten & gering \\
    \bottomrule
  \end{tabular}
  \caption{Übersicht der Forschungsmethoden}
\end{table}

\end{frame}

%%%%%%%%%%%%%%%%%%%%%%%%%%%%%%%%%%%%%%%%%%%%%%%%%%%%%%%%%%%%%%%%%%%%%%%%%%%%%%%
\begin{frame}{Zwei Spalten mit Abbildung}

\begin{columns}[T]
  \begin{column}{0.55\textwidth}
    \begin{itemize}
      \item Das NXT-Thema verwendet das charakteristische
        \emph{NXT-Grün} (RGB 0, 139, 146).
      \item Schriftart: \textbf{Segoe UI} (LuaLaTeX).
      \item Mathematik: Euler-Math.
      \item Die NXT-Logos stecken im NXT Logo Font und sind immer
        verfügbar; Hochschullogo und Skyline werden mit
        \texttt{\textbackslash IfFileExists} abgesichert.
    \end{itemize}
  \end{column}
  \begin{column}{0.4\textwidth}
    \begin{center}
      \nxtlogo[width=0.7\textwidth]\\[10pt]
      {\color{nxtgreen}\nxtlogofont\LARGE N X T}
    \end{center}
  \end{column}
\end{columns}

\end{frame}

%%%%%%%%%%%%%%%%%%%%%%%%%%%%%%%%%%%%%%%%%%%%%%%%%%%%%%%%%%%%%%%%%%%%%%%%%%%%%%%
\begin{frame}{Mathematik}

Die Euler'sche Identität verbindet fünf fundamentale Konstanten:
\[
  e^{i\pi} + 1 = 0
\]

\bigskip

Das Integral der Gaußschen Glockenkurve:
\[
  \int_{-\infty}^{+\infty} e^{-x^2}\,\mathrm{d}x = \sqrt{\pi}
\]

\bigskip

\begin{block}{Hinweis}
  Das Thema lädt \texttt{unicode-math} mit dem Font
  \texttt{Euler-Math} für typographisch ansprechenden Formelsatz.
\end{block}

\end{frame}

%%%%%%%%%%%%%%%%%%%%%%%%%%%%%%%%%%%%%%%%%%%%%%%%%%%%%%%%%%%%%%%%%%%%%%%%%%%%%%%
\begin{frame}{Farben – Grüntöne und Grautöne}

\begin{columns}[T]
  \begin{column}{0.48\textwidth}
    \textbf{NXT-Grüntöne und Akzente}\\[4pt]
    \cs{nxtgreen}\\[3pt]
    \cs{nxtgreen.2}\\[3pt]
    \cs{nxtgreen.3}\\[3pt]
    \cs{nxtgreen.4}\\[3pt]
    \cs{nxt.violet}\\[3pt]
    \cs{nxt.blue}\\[3pt]
    \cs{nxt.orange}\\[6pt]
    {\footnotesize\texttt{nxtcolor} ist ein Alias für \texttt{nxtgreen}.}
  \end{column}
  \begin{column}{0.48\textwidth}
    \textbf{Grautöne}\\[4pt]
    \cs{nxt.verylightgray}\\[3pt]
    \cs{nxt.gray0}\\[3pt]
    \cs{nxt.gray1}\\[3pt]
    \cs{nxt.gray3}\\[3pt]
    \cs{nxt.gray5}\\[3pt]
    \cs{nxt.gray7}\\[3pt]
    \cs{nxt.gray8}
  \end{column}
\end{columns}

\end{frame}

%%%%%%%%%%%%%%%%%%%%%%%%%%%%%%%%%%%%%%%%%%%%%%%%%%%%%%%%%%%%%%%%%%%%%%%%%%%%%%%
\begin{frame}{Farben – weitere Farben}

\begin{columns}[T]
  \begin{column}{0.48\textwidth}
    \cs{nxt.yellow}\\[3pt]
    \cs{nxt.red}\\[3pt]
    \cs{nxt.magenta}\\[3pt]
    \cs{nxt.violet}\\[3pt]
    \cs{footgray}
  \end{column}
  \begin{column}{0.48\textwidth}
    \cs{nxt.navy}\\[3pt]
    \cs{nxt.cyan}\\[3pt]
    \cs{nxt.olive}\\[3pt]
    \cs{nxt.lightolive}\\[3pt]
    \cs{nxt.lightblue}\\[3pt]
    \cs{nxt.darkblue}
  \end{column}
\end{columns}

\end{frame}

%%%%%%%%%%%%%%%%%%%%%%%%%%%%%%%%%%%%%%%%%%%%%%%%%%%%%%%%%%%%%%%%%%%%%%%%%%%%%%%
\begin{frame}[fragile]{Logos}

Die NXT-Logos kommen aus dem NXT Logo Font, der dem Paket beiliegt --
nichts herunterzuladen. Das optionale Argument nimmt \texttt{width},
\texttt{height} und das Farbschema:
\medskip

\begin{tabular}{@{}ll@{}}
  \verb|\nxtlogo|         & \nxtlogo[height=6mm]     \\[4pt]
  \verb|\nxtlogonut|      & \nxtlogonut[height=6mm]  \\[4pt]
  \verb|\nxthsrtlogo|     & \nxthsrtlogo[height=6mm] \\[4pt]
  \verb|\nxtlogostandard| & \nxtlogostandard[height=6mm] \\[4pt]
  \verb|schema=schwarz|   & \nxthsrtlogo[height=6mm, schema=schwarz] \\[4pt]
  \verb|schema=weisseZeile| &
    \colorbox{nxt.gray8}{\nxthsrtlogo[height=6mm, schema=weisseZeile]}
\end{tabular}

\end{frame}

%%%%%%%%%%%%%%%%%%%%%%%%%%%%%%%%%%%%%%%%%%%%%%%%%%%%%%%%%%%%%%%%%%%%%%%%%%%%%%%
\begin{frame}[fragile]{Hochschullogo und Skyline}

Hochschullogo und Skyline gehören der Hochschule und liegen dem Paket
nicht bei. \texttt{getnxtlogos} lädt sie von der Hochschul-Website und
installiert sie in den lokalen \texttt{texmf}-Baum:

\begin{verbatim}
sh getnxtlogos              # macOS, Linux
powershell -ExecutionPolicy Bypass
           -File getnxtlogos.ps1   # Windows
\end{verbatim}

\begin{tabular}{@{}ll@{}}
  \verb|\hsrtlogograu| & \hsrtlogograu[height=6mm] \\[4pt]
  \verb|\hsrtskyline|  & \hsrtskyline[height=4mm]
\end{tabular}
\medskip

Fehlen die Dateien, übergehen beide Makros sie stillschweigend.

\end{frame}

\end{document}
